\scanpage{12}
\mspage{3}
c) Soient $\unsure{\mathscr P|P},\mathscr P'$ deux catégories de Picard, épinglées par $M,N$\unsure{.|!} Pour qu'ils soient \unsure{équivalentes|équivalents} (avec leur épinglage) il f. et s. que les applications $\sigma_{\mathscr P},\sigma_{\mathscr P'}\colon M\longrightarrow {}_2N$ soient les \unsure{mêmes|$\overline{\mathrm{m}}$s|$=$s}. La catégorie des équivalences est alors \deleted{une}\inserted{un 1-torseur sous une} catégorie de Picard \uline{stricte}
% slanted
\inserted{d'invariants $\operatorname{Ext}^{1}(M,N)$ et $\operatorname{Hom}(M,N)$ \unsure{(définie par|(obtenus par} $R\underline{\operatorname{Hom}}(M,N)$ tronqué\unsure{\ldots})}

d) La classification à équivalence près des Cat. de Picard épinglées par $M,N$ est (comme groupe) $\simeq\operatorname{Hom}(M,{}_2N)$.

e) Revenant à la situation de c) et utilisant la multiplication de Baer des Pic cat. épinglées par $(M,N)$, on trouve que \deleted{\unsure{celles}} celles qui correspondent à un $\sigma\colon M\longrightarrow{}_2N$ donné forment une 2-catégorie qui est un \unsure{2|1}-torseur sous la 2-catégorie de Picard stricte formée des Cat de Pic. \uline{strictes} épinglées par $M,N$ \deleted{\unsure{E}} [qui est un 2-groupoïde connexe : Tous les 0-objets sont 2-équivalents,\ldots] \deleted{\unsure{les}}

f) Les Pic-catégories forment une \uline{2-$\otimes$-catégorie \deleted{\unsure{de}} AUC} de façon évidente, \deleted{\unsure{Si $\mathscr P,\mathscr P'$ et $\mathscr P$}} (on s'envoie dans \inserted{(et $=$ \uline{sur} les objets et flèches)} la 1-$\otimes$-catégorie \inserted{AUC} sous-jacente à la catégorie \deleted{\unsure{des ob}}\inserted{$D$}
% archive sheet number: 11
